\section{Introduction}

Diffusion models reverse a noising process by repeatedly predicting a dense
noise field. This dense prediction is effective, but it is also monolithic:
there is no direct way to ask how much of the denoising update has been
explained, whether the update has an interpretable coarse-to-fine order, or how
to stop after a valid partial prediction. Ordered visual tokenizers and
diffusion-autoregressive methods suggest that generation can benefit from
structured intermediate representations. The question in this paper is
narrower: can the dense pixel-space noise prediction itself be factorized into
ordered denoising tokens?

\method{} rewrites the predicted noise field as
\begin{equation}
  \epshat^{(K)} = \sum_{k=1}^{K} S_k\left(z^{(k)}\right),
  \label{eq:full-factorization}
\end{equation}
where each $z^{(k)}$ is a denoising token and $S_k$ expands only
that token into one dense component. Prefixes are meaningful by construction:
\begin{align}
  \epshat^{(1:m)} &= \sum_{k=1}^{m} S_k\left(z^{(k)}\right), \\
  \hat{x}^{(m)}_0 &= \frac{\xt - \sigma_t \epshat^{(1:m)}}{\alpha_t}.
  \label{eq:prefix-denoise}
\end{align}

The central design constraint is that the synthesis operator must not become a
decoder. In the default \method{} model, $S_k$ is token-only, bias-free,
shallow, local, and condition-free. It does not receive the noised image,
timestep, class label, text prompt, previous tokens, skip features, or a learned
constant. This restriction forces the token itself to carry the sample-specific
negative-noise information and gives the zero-token test a concrete meaning:
$S_k(0)=0$ by construction.

The current evidence supports a scoped method claim. \method{} learns ordered
denoising components across eight datasets, including ImageNet-family 64x64 and
256x256 settings and non-classification FFHQ/AFHQv2 rows. Compared with an
endpoint-only factorized predictor, \method{} consistently improves
  prefix/path alignment, while matched objective, prediction, and synthesis
  ablations expose why ordering and restricted synthesis matter. Generated-sample
distribution metrics do not show a stable unconditional generation-quality win,
which clarifies the present contribution: \method{} makes dense noise
prediction prefix-controllable without relying on a final VAE-style decoder.

\paragraph{Contributions.}
\begin{enumerate}
  \item We formulate dense pixel-space diffusion noise prediction as an ordered
  sum of denoising-token components.
  \item We introduce a restricted, condition-free synthesis operator whose
  token-only interface and bias-free construction prevent decoder-style
  shortcuts and make $S_k(0)=0$ testable.
  \item We train \method{} with a denoise-path objective that makes token
  prefixes approximate a partial denoising trajectory rather than only the final
  clean image.
  \item We evaluate prefix controllability, component ordering,
  non-degeneration, endpoint quality, matched short-run generation metrics, and
  tokenizer-related reconstruction baselines across the current eight-dataset
  evidence matrix.
\end{enumerate}
